\documentclass[book.tex]{subfiles}
\begin{document}

\chapter{Solving High Dimensional Parabolic PDEs }
\label{chap:high_dimensional}
\noindent\rule{11cm}{0.4pt} \\
\textbf{Prerequisites:} \autoref{chap:pinns} \\
\textbf{Difficulty Level:} ** \\
\noindent\rule{11cm}{0.4pt}
\vspace{5mm}

The Feynman-Kac formula provides a bridge between solutions to linear parabolic partial differential equations and stochastic differential equations. Recent work has extended the Feynman-Kac formula to broader classes of nonlinear parabolic differential equations. In particular, by leveraging the universal approximation capability of neural networks to model the gradient of the unknown solution, we can numerically solve high dimensional parabolic partial differential equations \cite{han2018solving}.

%"""
%Developing algorithms for solving high-dimensional partial differential equations (PDEs) has been an exceedingly difficult task for a long time, due to the notoriously difficult problem known as the “curse of dimensionality.” This paper introduces a deep learning-based approach that can handle general high dimensional parabolic PDEs. To this end, the PDEs are reformulated using backward stochastic differential equations and the gradient of the unknown solution is approximated by neural networks, very much in the spirit of deep reinforcement learning with the gradient acting as the policy function. Numerical results on examples including the nonlinear Black–Scholes equation, the Hamilton–Jacobi–Bellman equation, and the Allen–Cahn equation suggest that the proposed algorithm is quite effective in high dimensions, in terms of both accuracy and cost. This opens up possibilities in economics, finance, operational research, and physics, by considering all participating agents, assets, resources, or particles together at the same time, instead of making ad hoc assumptions on their interrelationships.
%"""


\section{Mathematical Formulation} 

Consider the family of semilinear parabolic PDEs given by

\begin{align}
    \frac{\partial u}{\partial t}(t, x) + \frac{1}{2}\mathrm{Tr}\left ( \sigma \sigma^T (t, x)(\mathrm{Hess}_x u)(t, x) \right ) + \nabla u(t, x) \cdot \mu(t, x) + f \left (t, x, u(t, x), \sigma^T(t, x) \nabla u(t, x) \right ) = 0 
\end{align}

Here $t$ is the time variable, $x$ is a $d$-dimensional space variable.

\begin{align} 
u(T, x) = g
\end{align}

$\mu$ is a vector valued variable and $\sigma$ is a $d \times d$ matrix valued function. 

\begin{align}
    &\mu: \mathbb{R} \to \mathbb{R}^d \\
    &\sigma: \mathbb{R} \to \mathbb{R}^{d \times d}
\end{align}

Suppose also that we have terminal condition at time $T$

Let $W_s$ be Brownian motion and let $X_t$ be a stochastic process that satisfies the equation

\begin{align}
    X_t &= \xi + \int_0^t \mu(s, X_s) ds + \int_0^t \sigma(s, X_s) dW_s
\end{align}

We can then express the solution to our original PDE as the following stochastic differential equation (SDE) by the Feynman-Kac formula

\begin{align}
u(t, X_t) - u(0, X_0) &= - \int_0^t f(s, X_s, u(s, X_s), \sigma(s, X_s) \nabla u (s, X_s)) ds + \int_0^t [\nabla u(s, X_s)]^T \sigma(s, X_s) dW_s
\end{align}

In order to turn this integral into a numerical algorithm, we partition $[0, 1]$ into $t_1,\dotsc, t_n$. We can then write down the Euler approximation

\begin{align}
    X_{t_{n+1}} - X_{t_n} \approx \mu(t_n, X_{t_n}) \Delta t_n + \sigma(t_n, X_{t_n}) \Delta W_n
\end{align}

This equation can be used to sample a path $X_t$. We can then construct an approximate solution of the SDE as

\begin{align}
u(t_{n+1}, X_{t_{n+1}}) - u(t_n, X_{t_n}) \approx - f\left (t_n, X_{t_n}, u(t_n, X_{t_n}), \sigma(t_n, X_{t_n}) \nabla u(t_n, X_{t_n}) \right) \Delta t_n + [\nabla u(t_n, X_{t_n})]^T \sigma(t_n, X_{t_n}) \Delta W_n
\end{align}

The main challenge of this equation though is that the update is formulated in terms of $\nabla u$, where $u$ is itself unknown. In order to use this equation, we have to find a way to compute this quantity. The idea is that we can approximate this gradient quantity by a neural network

\begin{align}
    x \mapsto \sigma^T(t_n, X_{t_n})\nabla u(t_n, X_{t_n}) \approx g_n(t_n, X_{t_n}, \theta)
\end{align}

where $\theta$ is a set of learnable parameters. Using this approximation we can compute an approximate solution $\hat{u}$ parameterized by $\theta$ given a sampled path $X_t$. To train parameters $\theta$, we can use the fact that we know terminal condition $g$ to impose a loss

\begin{align}
\mathcal{L}(\theta) &= \mathbb{E}_{X_t} \left [ \left | g(X_{t_n}) - \hat{u}(X_t, W_t)\right |^2 \right ]
\end{align}

This loss can be optimized to find $\theta$ that provides an approximate solution to the original PDE.

\section{Network Architecture}

The gradient term $\nabla u$ is approximated using a multilayer feed forward network that maps $X_{t_n}$ to the estimated term

\begin{align}
    X_{t_n} \to h^1_n \to h^2_n \dotsc \to h^H_n \to \nabla u(t_n, X_{t_n})
\end{align}
We can convert this definition into Physika code

\begin{lstlisting}[language=python]
class SpatialGradient($\theta$, $n_{\textrm{layers}}$):

  def $\lambda$(X):
    return fcnet(X, $\theta, n_{\textrm{layers}})$

\end{lstlisting}

The forward iteration 
\begin{align}
u(t_n, X_{t_n}) \mapsto u(t_{n+1}, X_{t_{n+1}})
\end{align}
can be computed directly from the definition earlier and the known functions $f$ and $\sigma$.
\begin{align}
    (u(t_n, X_{t_n}), \nabla u(t_n, X_{t_n}), W_{t_{n+1}} %- W_{t_n}) \mapsto u(t_{n+1}, X_{t_{n+1}})
\end{align}

\begin{lstlisting}[language=python]

class ForwardIteration($\theta$, $n_{\textrm{layers}}$):

  def $\lambda$(X):
    return fcnet(X, $\theta, $n_{\textrm{layers}}$)$

\end{lstlisting}


This network is a shortcut connecting blocks at different times. 
\begin{align}
    X_{t_n}, W_{t_{n+1}} - W_{t_n} \mapsto X_{t_{n+1}}
\end{align}

\begin{lstlisting}[language=python]

class TimeUpdate($\theta$, n_layers):

  def $\lambda$(X):
    return fcnet(X, $\theta, n_layers)$

\end{lstlisting}

One fact to note about these networks is that they can become very deep. Suppose that spatial gradient network has $H$ hidden layers. Then the total network has $(H+1)(N-1)$ layers in total for a discretization with $N$ points total.

%\textcolor{red}{TODO: The scheme seemms very general? Why does this only hold for semilinear parabolic equations.}

%\section{Parabolic PDEs}

%Suppose we have a second order PDE given by 
%\begin{align}
%    A u_{xx} + 2 B u_{xy} + C u_{yy} + D u_x + E u_y + F = 0
%\end{align}
%We say that this PDE is parabolic if $B^2 - AC = 0$. \textcolor{red}{TODO: I don't have any intuition for this definition. Why should this PDE being parabolic make a difference?}

%Monte Carlo methods can efficiently solve high dimensional linear parabolic PDEs \textcolor{red}{TODO:CITE}. Solving semilinear parabolic PDEs is harder.
\end{document}